\documentclass[a4paper,12pt]{article}
\usepackage[T2A]{fontenc}
\usepackage[utf8]{inputenc}
\usepackage[russian]{babel}
\usepackage{amsmath,amssymb,mathtools}
\usepackage{geometry}
\geometry{left=22mm,right=22mm,top=23mm,bottom=23mm}
\usepackage{microtype}
\usepackage{tikz}
\usetikzlibrary{calc,arrows.meta,positioning,shapes.geometric}
\usepackage{tabularx,booktabs}
\usepackage{enumitem}
\usepackage{fancyhdr}
\usepackage{hyperref}
\usepackage{xcolor}
\usepackage{parskip}
\definecolor{accent}{HTML}{008080}
\definecolor{darkaccent}{HTML}{005A5A}
\definecolor{textgray}{HTML}{333333}
\definecolor{subtle}{HTML}{777777}
\hypersetup{colorlinks=true,linkcolor=darkaccent,urlcolor=darkaccent}
\setlength{\parindent}{0pt}
\setlength{\parskip}{5pt plus 1pt}
\setlist[enumerate]{leftmargin=*,itemsep=6pt,topsep=5pt}
\setlist[itemize]{leftmargin=*,itemsep=3pt,topsep=3pt}
\pagestyle{fancy}
\fancyhf{}
\fancyhead[L]{\footnotesize Логарифмы и логарифмические уравнения}
\fancyhead[R]{\footnotesize Даниил Хрястов}
\renewcommand{\headrulewidth}{0.4pt}
\renewcommand{\headrule}{\hbox to\headwidth{\color{accent}\leaders\hrule height \headrulewidth\hfill}}
\fancyfoot[C]{\footnotesize\thepage}
\newcommand{\lessonsection}[1]{\section*{\color{darkaccent}#1}\vspace{-5pt}{\color{accent}\rule{\linewidth}{0.45pt}}\vspace{5pt}}
\newcommand{\smallhead}[1]{\vspace{4pt}{\color{darkaccent}\bfseries #1}\par}
\newcommand{\mathnote}[1]{\textcolor{darkaccent}{\textbf{#1}}}
\newcommand{\exercise}[2]{\item \textbf{#1}\\[-2pt] \[#2\]}
\setlength{\headheight}{25pt}
\begin{document}

\begin{titlepage}
\thispagestyle{empty}
\begin{center}
\vspace*{31mm}
{\color{darkaccent}\normalsize МАТЕМАТИКА\par}
\vspace{19mm}
{\LARGE Чек-лист\par}
\vspace{14mm}
{\color{accent}\Large Логарифмы: преобразования\\[3pt] и решение уравнений\par}
\vspace{19mm}
{\normalsize Подготовка к ЕГЭ по профильной математике\par}
\vfill
{\normalsize Даниил Хрястов\par}
\vspace{12mm}
{\normalsize 13 ноября 2026 г.\par}
\vspace{21mm}
{\small Лёвин Артём Александрович\\[3pt]
эксклюзивно для Даниила Хрястова\par}
\vspace{20mm}
\begin{tikzpicture}
\draw[accent,line width=.9pt]
 (-4.5,0) .. controls (-3.1,.55) and (-1.8,-.47) .. (0,0)
 .. controls (1.6,.48) and (3.2,-.45) .. (4.5,0);
\end{tikzpicture}
\end{center}
\end{titlepage}
\setcounter{page}{1}

\lessonsection{1. Что означает логарифм}
Логарифм --- это показатель степени, в которую нужно возвести основание, чтобы получить данное число:
\[
\boxed{\log_a b=c\quad\Longleftrightarrow\quad a^c=b,\qquad a>0,\ a\ne1,\ b>0.}
\]
В записи $\log_a b$ число $a$ называется \emph{основанием}, а $b$ --- \emph{аргументом} логарифма (на занятиях также: «тело логарифма»).

\smallhead{Пример 1. Найдите $\log_2 64$.}
\[
64=2^6\quad\Longrightarrow\quad\log_2 64=\log_2 2^6=6.
\]
Здесь полезно сразу представить $64$ степенью двойки.

\smallhead{Основное логарифмическое тождество}
\[
\boxed{a^{\log_a b}=b,\qquad a>0,\ a\ne1,\ b>0.}
\]
Именно эту формулу называют \emph{основным логарифмическим тождеством}. Например,
$5^{\log_5 7}=7$.

\smallhead{Пример 2. Степень степени.}
\[
\left(3^{\log_3 2}\right)^{\log_2 3}
 =2^{\log_2 3}=3.
\]
Можно использовать и свойство $(a^m)^n=a^{mn}$. Главное --- определить, какие основания и показатели уже удобно сопоставить.

\lessonsection{2. Формулы для преобразования логарифмов}
Во всех формулах ниже $a>0$, $a\ne1$; аргументы логарифмов положительны.
\[
\begin{aligned}
\log_a(xy)&=\log_a x+\log_a y, &&x>0,\ y>0;\\[4pt]
\log_a\!\left(\frac{x}{y}\right)&=\log_a x-\log_a y, &&x>0,\ y>0;\\[4pt]
\log_a(x^r)&=r\log_a x, &&x>0;\\[4pt]
\log_a a&=1,\qquad \log_a 1=0.&&
\end{aligned}
\]
\smallhead{Смена основания}
\[
\boxed{\log_a b=\frac{\log_c b}{\log_c a}},\quad
\boxed{\log_a b=\frac{1}{\log_b a}}.
\]
Во второй формуле $a,b>0$ и $a,b\ne1$. В первой формуле $c>0$, $c\ne1$.

\textcolor{darkaccent}{\textbf{Обратите внимание:}} знак степени можно вынести перед логарифмом, когда в степень возведён \emph{весь аргумент}: $\log_a(b^3)=3\log_a b$. Для $\log_a(b^3/a)$ сперва учитывают деление внутри аргумента.

\lessonsection{3. Разбираем выражения по действиям}
\smallhead{Пример 3. Дано $\log_b a=\frac17$. Вычислите $\log_a(ab^3)$.}
\[
\begin{aligned}
\log_a(ab^3)&=\log_a a+\log_a b^3\\
&=1+3\log_a b\\
&=1+3\cdot\frac{1}{\log_b a}\\
&=1+3\cdot7=\boxed{22}.
\end{aligned}
\]
Сначала раскладываем произведение на сумму логарифмов, затем выносим показатель и применяем формулу взаимного перехода между основаниями.

\smallhead{Пример 4. Дано $\log_a b=5$. Вычислите $-\log_a\!\left(\dfrac{b^3}{a}\right)$.}
\[
\begin{aligned}
-\log_a\!\left(\frac{b^3}{a}\right)
&=-\bigl(\log_a b^3-\log_a a\bigr)\\
&=-\bigl(3\log_a b-1\bigr)\\
&=-\bigl(3\cdot5-1\bigr)=\boxed{-14}.
\end{aligned}
\]
\textcolor{darkaccent}{\textbf{Правило записи:}} при подстановке целого выражения ставьте скобки. Иначе знак «минус» перед логарифмом может изменить только первое слагаемое и привести к ошибке.

\lessonsection{4. Первый тип уравнений: логарифм равен числу}
Из определения логарифма следует:
\[
\boxed{\log_a F(x)=c\quad\Longleftrightarrow\quad F(x)=a^c},
\]
причём исходный логарифм должен быть определён. На занятии этот переход запоминали как «\emph{циклопирование}».

\smallhead{Пример 5. Постоянное основание.}
\[
\begin{cases}
\log_2 x=5,\\x>0
\end{cases}
\Longrightarrow x=2^5=\boxed{32}.
\]

\smallhead{Пример 6. Переменная в основании.}
\[
\begin{cases}
\log_x16=2,\\x>0,\quad x\ne1
\end{cases}
\Longrightarrow x^2=16 \Longrightarrow x=\pm4.
\]
Из двух кандидатов $-4$ нарушает условие $x>0$. \mathnote{Ответ:} $x=4$.

\smallhead{Пример 7. Отрицательный знак внутри основания.}
\[
\begin{cases}
\log_{-x}16=2,\\ -x>0,\quad -x\ne1
\end{cases}
\Longrightarrow (-x)^2=16\Longrightarrow x=\pm4.
\]
Здесь требуется $x<0$ и $x\ne-1$. \mathnote{Ответ:} $x=-4$.

\textcolor{darkaccent}{\textbf{Вывод:}} отрицательный кандидат не исключается автоматически. Всё решают ограничения \emph{конкретного} логарифма.

\lessonsection{5. Второй тип: логарифм равен логарифму}
Если основания одинаковые и перед каждым логарифмом нет дополнительных множителей или слагаемых, их аргументы можно приравнять:
\[
\boxed{\log_a F(x)=\log_a G(x)\iff F(x)=G(x)}
\]
при $a>0$, $a\ne1$ и условиях $F(x)>0$, $G(x)>0$.

\smallhead{Пример 8. Разные основания: сначала преобразуем.}
\[
\log_2 x=\log_4 5=\frac{\log_2 5}{\log_2 4}
=\frac12\log_2 5=\log_2\sqrt5.
\]
С учётом $x>0$ получаем $\boxed{x=\sqrt5}$.

\smallhead{Пример 9. Слагаемое вне логарифма.}
\[
\begin{aligned}
\log_2 x&=\log_2 5+1
=\log_2 5+\log_2 2\\
&=\log_2(5\cdot2)=\log_2 10
\quad\Longrightarrow\quad\boxed{x=10}.
\end{aligned}
\]
Здесь единица находится \emph{вне} логарифма. Если бы она была внутри, запись выглядела бы как $\log_2(5+1)$, и результат был бы другим.

\smallhead{Пример 10. Приводим основание к степени двойки.}
\[
\log_8\!\left(2^{8x-4}\right)=4,
\qquad 8=2^3,
\]
\[
\frac{8x-4}{3}=4
\quad\Longrightarrow\quad8x-4=12
\quad\Longrightarrow\quad\boxed{x=2}.
\]
Перед преобразованием полезно выяснить, можно ли упростить степень или аргумент с помощью уже известной формулы.

\lessonsection{6. Проверка ограничений и сложное уравнение}
\smallhead{Пример 11. Смена основания и отбор корней.}
Решим уравнение
\[
\log_2\!\left(\frac{x}{6}\right)=\log_{1/2}(x+1).
\]
\textbf{Шаг 1. ОДЗ:} $x/6>0$ и $x+1>0$, то есть $x>0$.

\textbf{Шаг 2. Сменим основание:}
\[
\log_{1/2}(x+1)=-\log_2(x+1)
=\log_2\!\left(\frac{1}{x+1}\right).
\]
\textbf{Шаг 3. Приравняем аргументы и решим:}
\[
\frac{x}{6}=\frac{1}{x+1}
\quad\Longleftrightarrow\quad x(x+1)=6
\quad\Longleftrightarrow\quad x^2+x-6=0.
\]
\[
(x+3)(x-2)=0\quad\Longrightarrow\quad x=-3\ \text{или}\ x=2.
\]
\textbf{Шаг 4. Отбор:} $x=-3$ не входит в ОДЗ. \mathnote{Ответ:} $\boxed{x=2}$.

\smallhead{Полезные формулы для поиска неизвестного}
С этими преобразованиями важно работать без ошибок знака:
\[
\boxed{A-B=C\ \Longrightarrow\ B=A-C},
\qquad
\boxed{\frac{A}{B}=C\ \Longrightarrow\ B=\frac AC}\quad (B,C\ne0).
\]
Например, $4-x=128\Longrightarrow x=4-128=-124$.

\textcolor{darkaccent}{\textbf{Правило оформления:}} показывайте существенные переходы на бумаге. Ошибочную строку достаточно зачеркнуть одной косой чертой так, чтобы она оставалась читаемой при проверке.

\clearpage
\thispagestyle{fancy}
\begin{center}
{\color{darkaccent}\Large Алгоритм: решение логарифмического уравнения}\par
\vspace{7mm}
\end{center}
\begin{center}
\begin{tikzpicture}[
>=Stealth, node distance=9mm and 13mm,
terminal/.style={ellipse,draw=darkaccent,very thick,align=center,minimum width=39mm,minimum height=9mm,font=\small},
action/.style={rectangle,draw=accent,rounded corners=1pt,align=center,minimum width=51mm,text width=49mm,minimum height=10mm,inner sep=3pt,font=\small},
condition/.style={diamond,aspect=2.6,draw=darkaccent,align=center,text width=37mm,inner sep=2pt,font=\small},
arrow/.style={->,line width=.75pt,draw=darkaccent},
lab/.style={font=\scriptsize,fill=white,inner sep=1.3pt}]
\node[terminal] (start) at (0,0) {Начало};
\node[action] (odz) at (0,-1.55) {Запишите ОДЗ исходных\ логарифмов};
\node[action] (prep) at (0,-3.15) {Приведите уравнение к простому виду: преобразуйте суммы, степени и основания};
\node[condition] (single) at (0,-5.5) {Один логарифм\ равен числу?};
\node[action] (exp) at (-4.5,-8.0) {$\log_a F(x)=c$\\$F(x)=a^c$};
\node[condition] (double) at (3.8,-8.0) {Два логарифма\ одного основания?};
\node[action] (eq) at (3.8,-10.15) {$\log_a F(x)=\log_a G(x)$\\$F(x)=G(x)$};
\node[action] (transform) at (6.35,-5.4) {Преобразуйте уравнение, затем заново определите его вид};
\node[action] (solve) at (0,-12.25) {Решите полученное\ алгебраическое уравнение};
\node[action] (filter) at (0,-14.05) {Проверьте все найденные\ значения по ОДЗ};
\node[terminal] (end) at (0,-15.6) {Ответ};
\draw[arrow] (start)--(odz);
\draw[arrow] (odz)--(prep);
\draw[arrow] (prep)--(single);
\draw[arrow] (single.west) -| node[lab,near start,above] {да} (exp.north);
\draw[arrow] (single.east) -| node[lab,near start,above] {нет} (double.north);
\draw[arrow] (double)--node[lab,right] {да}(eq);
\draw[arrow] (double.east) -| node[lab,near start,above] {нет} (transform.south);
\draw[arrow] (transform.north) |- (prep.east);
\draw[arrow] (exp.south) |- (solve.west);
\draw[arrow] (eq.south) |- (solve.east);
\draw[arrow] (solve)--(filter);
\draw[arrow] (filter)--(end);
\end{tikzpicture}
\end{center}
\clearpage

\lessonsection{7. Частые ошибки и контрольный список}
\begin{itemize}
\item $\log_a(x+y)\ne\log_a x+\log_a y$ в общем случае. Формула суммы логарифмов относится к \emph{произведению} аргументов.
\item $\log_a(x/y)=\log_a x-\log_a y$ требует $x>0$, $y>0$.
\item В равенстве $\log_a F(x)=\log_b G(x)$ сначала сравните основания. При $a\ne b$ аргументы сразу приравнивать нельзя.
\item Различайте $\log_2 5+1$ и $\log_2(5+1)$: скобки определяют границу аргумента.
\item Записывайте ОДЗ \emph{до} преобразований, а затем проверяйте каждый найденный корень.
\item Не теряйте минус при раскрытии скобок: $-(u-v)=-u+v$.
\item При переносе слагаемого или поиске вычитаемого полезно подставить результат обратно.
\end{itemize}

\smallhead{Перед сдачей решения проверьте себя}
\begin{enumerate}[label={\(\square\) \arabic*.},itemsep=3pt]
\item Я записал(а) ограничения для всех логарифмов.
\item Я показал(а) основные преобразования и аккуратно расставил(а) скобки.
\item Если приравниваю аргументы, основания одинаковы и перед логарифмами нет лишних действий.
\item Я решил(а) полученное уравнение и проверил(а) его корни по ОДЗ.
\item Я записал(а) только допустимые решения.
\end{enumerate}

\clearpage
\lessonsection{Тренировочный блок}
{\color{subtle}\small Путь А. 10 заданий: от основных формул к составным логарифмическим уравнениям. Решения оформляйте последовательно и указывайте ОДЗ, если есть неизвестная под знаком логарифма или в его основании.}\par
\vspace{4pt}
\smallhead{Базовый уровень}
\begin{enumerate}[label=\textbf{\arabic*.},itemsep=8pt]
\item Вычислите $\displaystyle\log_2 128$.
\item Вычислите $\displaystyle 7^{\log_7 11}$.
\item Известно, что $\displaystyle\log_b a=\frac15$. Вычислите $\displaystyle\log_a(ab^2)$.
\item Известно, что $\displaystyle\log_a b=4$. Вычислите $\displaystyle-\log_a\!\left(\frac{b^3}{a}\right)$.
\end{enumerate}
\smallhead{Средний уровень}
\begin{enumerate}[label=\textbf{\arabic*.},start=5,itemsep=8pt]
\item Решите уравнение $\displaystyle\log_3(x-1)=2$.
\item Решите уравнение $\displaystyle\log_x81=4$.
\item Решите уравнение $\displaystyle\log_{-x}25=2$.
\item Решите уравнение $\displaystyle\log_2(7-x)=6$.
\end{enumerate}
\smallhead{Повышенный уровень}
\begin{enumerate}[label=\textbf{\arabic*.},start=9,itemsep=8pt]
\item Решите уравнение $\displaystyle\log_2 x=\log_4 5+1$.
\item Решите уравнение $\displaystyle\log_2\!\left(\frac{x}{6}\right)=\log_{1/2}(x+1)$.
\end{enumerate}

\clearpage
\lessonsection{Ответы к тренировочному блоку}
\renewcommand{\arraystretch}{1.5}
\begin{center}
\begin{tabularx}{\linewidth}{@{}>{\raggedleft\arraybackslash}p{17mm} X @{} }
\toprule
\textbf{№} & \textbf{Ответ} \\
\midrule
1 & $7$ \\
2 & $11$ \\
3 & $11$ \\
4 & $-11$ \\
5 & $x=10$ \\
6 & $x=3$ \\
7 & $x=-5$ \\
8 & $x=-57$ \\
9 & $x=2\sqrt5$ \\
10 & $x=2$ \\
\bottomrule
\end{tabularx}
\end{center}
\vspace{11mm}
\smallhead{Итоговая памятка}
\begin{tabularx}{\linewidth}{@{}p{48mm}X@{}}
\toprule
\textbf{Сигнал в задании} & \textbf{Что сделать}\\
\midrule
Логарифм равен числу & Перейти к степени основания; проверить ОДЗ.\\
Два логарифма & Привести основания к одному; убрать внешние действия.\\
Произведение или дробь & Использовать сумму или разность логарифмов.\\
Степень внутри логарифма & Проверить аргумент и вынести допустимый показатель.\\
Разные основания & Применить формулу смены основания.\\
Несколько корней & Оставить только удовлетворяющие исходной ОДЗ.\\
\bottomrule
\end{tabularx}
\vfill
\begin{center}
{\footnotesize\color{subtle} Лёвин Артём Александрович --- эксклюзивно для Даниила Хрястова}
\end{center}
\end{document}
